\documentclass{article}

\usepackage[utf8]{inputenc}
\usepackage[T1]{fontenc}
\usepackage{helvet}
\usepackage{geometry}
\usepackage{xcolor}
\usepackage{eso-pic}
\usepackage{fancyhdr}
\usepackage{parskip}
\usepackage{microtype}
\usepackage[english, ukrainian]{babel}
\usepackage{url}
\usepackage{amsmath}
\usepackage{amssymb}
\usepackage{longtable}
\usepackage{booktabs}
\usepackage{hyperref}
\usepackage{titlesec}

\definecolor{nato}{HTML}{293757}
\definecolor{footergray}{gray}{0.65}

\geometry{a4paper,left=3cm,right=3cm,top=3cm,bottom=3cm}
\renewcommand{\familydefault}{\sfdefault}
\setcounter{secnumdepth}{3}

\titleformat{\section}{\color{nato}\Huge\bfseries}{\thesection.}{1em}{}
\titlespacing*{\section}{0pt}{30pt}{12pt}
\titleformat{\subsection}{\color{nato}\LARGE\bfseries}{\thesubsection.}{1em}{}
\titlespacing*{\subsection}{0pt}{20pt}{8pt}
\titleformat{\subsubsection}{\color{nato}\Large\bfseries}{\thesubsubsection.}{1em}{}
\titlespacing*{\subsubsection}{0pt}{10pt}{6pt}

\fancyhf{}
\fancyhead[R]{\small\color{footergray} 29 червня 2026}
\fancyfoot[L]{\small\color{footergray} Технічні вимоги ERP/1. Модуль IAS}
\fancyfoot[R]{\small\color{footergray}\thepage}
\newcommand\HUGE[1]{\fontsize{40}{40}\selectfont #1}
\renewcommand{\headrulewidth}{0pt}
\renewcommand{\footrulewidth}{0.4pt}
\renewcommand{\footrule}{\color{footergray}\hrule width \textwidth height 0.4pt}

\begin{document}
\pagestyle{fancy}

\begin{titlepage}
  \AddToShipoutPictureBG*{\AtPageLowerLeft{\color{nato}\rule{\paperwidth}{\paperheight}}}
  \color{white}\thispagestyle{empty}\vspace*{4cm}\centering
  {\Huge \textbf{ERP/1: Авторизація} \par} \vspace{2cm}
  {\HUGE \textbf{Технічні вимоги} \par} \vspace{2.5cm}
  {\LARGE \textbf{до системи керування цифровими ідентичностями, сертифікатами та політиками безпеки} \par} \vspace{2.5cm}
  \vfill
  {\large 2026 \copyright\ Системи Електронного Урядування \par}
\end{titlepage}

\newpage
\begin{center}{\Huge\color{nato}\bfseries Зміст\par}\end{center}
\vspace{40pt}
\tableofcontents

\begin{abstract}
ERP/1: Служба ідентифікації та авторизації (IAS) є центральним керуючим компонентом безпеки Zencrypted-середовища на базі Erlang/OTP. Поточна реалізація веде стійкий реєстр доменних об'єктів у KVS/Mnesia, підтримує граф зв'язків користувач--пристрій--сертифікат--профіль безпеки--VPN-сервіс, перевірку CSR і X.509, покрокове налаштування пристрою, доставку бажаного стану до VPN через Erlang RPC, ревізії, аудит доставки, звіряння стану та відновлення після збоїв. Інтеграція з CA у поточній версії містить перехідний CMP-шлях через зовнішній командний інтерфейс OpenSSL; цільовою є повне усунення запуску OpenSSL із середовища виконання IAS та взаємодія з CA через керований клієнт або сервісний адаптер. Вимоги щодо BPE/BPMN, EST/CMC, CRL/OCSP, КСЗІ, КЕП та національних криптографічних алгоритмів у цьому документі позначаються окремо як цільові або нормативні вимоги і не трактуються як уже реалізовані функції.
\end{abstract}

\newpage
\section{Умовні скорочення та визначення}
\begin{longtable}{p{4cm}p{10cm}}
\toprule
\textbf{Термін / Скорочення} & \textbf{Значення} \\
\midrule
IAS & Identity and Authorization Service \\
CA & Certificate Authority \\
VPN & Virtual Private Network \\
PKI & Public Key Infrastructure \\
CSR & Certificate Signing Request \\
CMP & Certificate Management Protocol (RFC 4210) \\
EST & Enrollment over Secure Transport (RFC 7030) \\
CMC & Certificate Management over CMS (RFC 5272) \\
BPE & Business Process Engine \\
BPMN & Business Process Model and Notation \\
KVS & Key-Value Storage \\
ABAC & Attribute-Based Access Control \\
FSM & Finite State Machine \\
OTP & Open Telecom Platform \\
\bottomrule
\end{longtable}

\section{Статуси вимог і тверджень}
Щоб не змішувати фактичну реалізацію, план розвитку та нормативні наміри, у документі використовуються такі статуси:
\begin{longtable}{p{3.2cm}p{10.8cm}}
\toprule
\textbf{Статус} & \textbf{Зміст} \\
\midrule
Реалізовано & Підтверджено поточним кодом і тестами IAS. \\
Частково реалізовано & Реалізовано частково або лише для конкретного доменного сценарію. \\
Цільовий стан & Цільова архітектура або функція наступної ітерації. \\
Нормативна вимога & Нормативна, сертифікаційна або організаційна вимога; не є доказом реалізації. \\
\bottomrule
\end{longtable}

\section{Загальні відомості}
\subsection{Призначення}
IAS централізує керування цифровими ідентичностями, пристроями, сертифікатами, профілями безпеки, VPN-сервісами та зв'язками між ними. IAS є джерелом бажаного стану для керованих VPN-конфігурацій і координує підготовку та передавання конфігурацій, але не бере участі у пересиланні VPN-трафіку.

\subsection{Поточний функціональний контур}
\begin{itemize}
  \item \textbf{Реалізовано:} стійке доменне сховище KVS/Mnesia та версіоновані записи.
  \item \textbf{Реалізовано:} граф користувач--пристрій--сертифікат--профіль безпеки--VPN-сервіс.
  \item \textbf{Реалізовано:} попередній перегляд та імпорт OVPN, перевірка CSR і X.509, ланцюжок походження сертифікатів та оцінювання довіри у межах поточного IAS/VPN домену.
  \item \textbf{Реалізовано:} майстер прив’язаного до пристрою налаштування VPN з перевірками переходів, готовності матеріалів, виправленням помилок та відновленням чернетки.
  \item \textbf{Реалізовано:} IAS $\rightarrow$ VPN Erlang RPC, канонічні команди, ревізії, контрольна сума, ідемпотентність, журналювання доставки і захист від застарілих ревізій.
  \item \textbf{Реалізовано:} звіряння стану, виявлення осиротілих записів, відновлення або виведення з експлуатації та відновлення після перезапуску.
  \item \textbf{Частково реалізовано:} випуск сертифікатів через CA; поточний CMP-шлях через зовнішній командний інтерфейс OpenSSL є перехідним, а OpenSSL має бути повністю усунений із середовища виконання IAS.
  \item \textbf{Цільовий стан:} керування повним життєвим циклом через \texttt{synrc/bpmn}/BPE.
\end{itemize}

\subsection{Нормативна рамка}
Застосовні закони, ДСТУ, ISO/IEC/IEEE, NIST та RFC розглядаються як джерела вимог і критеріїв аудиту. Факт згадки стандарту не означає, що сертифікація, атестація або повна технічна відповідність уже завершені.

\section{Архітектурні межі}
\subsection{IAS і CA}
\begin{itemize}
  \item \textbf{Поточний стан / Частково реалізовано:} IAS містить адаптер випуску сертифікатів за CMP, який запускає зовнішній командний інтерфейс OpenSSL і звертається до точки доступу CA для CMP.
  \item \textbf{Цільовий стан:} IAS не запускає OpenSSL або інші зовнішні криптографічні засоби командного рядка. CSR та контекст випуску передаються до CA через керований клієнт або сервісний адаптер; CMP може бути збережений як протокол взаємодії.
  \item Приватний ключ пристрою не повинен передаватися до IAS або CA.
  \item EST і CMC не входять до поточної реалізації IAS і можуть розглядатися лише як майбутні адаптери.
\end{itemize}

\subsection{IAS і VPN}
\begin{itemize}
  \item IAS є джерелом бажаного стану керованих VPN-вузлів і сервісів.
  \item Доставка виконується через розподілений Erlang RPC канонічними командами з ревізією та контрольною сумою.
  \item VPN застосовує команди і зберігає власний стійкий стан виконання.
  \item IAS веде журналювання доставки, стан очікування/застосування, виявляє застарілі ревізії та відновлює доставку після перезапуску або повторного з’єднання.
  \item Звіряння стану порівнює авторитетний стан IAS зі знімком стану VPN і формує інциденти та дії відновлення або виведення з експлуатації.
\end{itemize}

\subsection{Веб-інтерфейс і сховище}
\begin{itemize}
  \item Браузер $\leftrightarrow$ IAS: N2O/NITRO/Cowboy.
  \item Зберігання даних IAS: KVS/Mnesia через доменні модулі сховища.
  \item \texttt{ias\_demo\_store} використовується як фасад/проєкція для інтерфейсу користувача та демо-сценаріїв і не визначає окрему продуктивну БД.
\end{itemize}

\section{Функціональні вимоги}
\subsection{Доменна модель}
IAS повинен підтримувати щонайменше такі класи об'єктів:
\begin{itemize}
  \item користувач, пристрій, сертифікат, профіль безпеки, політика безпеки, VPN-сервіс;
  \item випуск і перевірка сертифікатів;
  \item транзакція підготовки конфігурації, стан доставки та журналювання доставки;
  \item граф зв’язків між користувачами, пристроями, сертифікатами, профілями та сервісами;
  \item Інцидент звіряння стану, відновлення або усунення осиротілого запису та бар’єр виведення з експлуатації.
\end{itemize}

\subsection{Життєвий цикл майстра налаштування VPN}
Центральний керований процес поточної версії має таку послідовність:
\begin{enumerate}
  \item Схема.
  \item Користувач.
  \item Пристрій.
  \item Профіль і політика безпеки.
  \item VPN-сервіс.
  \item Сертифікат CA.
  \item Клієнтський сертифікат.
  \item Зв’язки.
  \item Готовність матеріалів.
  \item Підготовка та передавання конфігурації.
\end{enumerate}

Кроки майстра та правила переходу наведено в таблиці~\ref{tab:wizard-steps}.

\begin{longtable}{p{2.6cm}p{4.0cm}p{4.2cm}p{3.0cm}}
\caption{Кроки майстра налаштування VPN}\label{tab:wizard-steps}\\
\toprule
\textbf{Крок} & \textbf{Необхідні матеріали} & \textbf{Умова переходу} & \textbf{Результат} \\
\midrule
\endfirsthead
\toprule
\textbf{Крок} & \textbf{Необхідні матеріали} & \textbf{Умова переходу} & \textbf{Результат} \\
\midrule
\endhead
Схема & Обрана схема налаштування & Схему вибрано та збережено в чернетці & Зафіксовано варіант подальшого процесу \\
Користувач & Наявний або створений користувач & Користувача вибрано та його ідентифікатор збережено & Зафіксовано власника пристрою \\
Пристрій & Наявний або створений пристрій & Пристрій належить вибраному користувачеві & Зафіксовано керований пристрій \\
Профіль і політика безпеки & Профіль безпеки та пов'язана політика & Профіль доступний для обраного сценарію & Зафіксовано правила перевірки й доступу \\
VPN-сервіс & Зареєстрований VPN-сервіс & Сервіс доступний і сумісний із профілем & Зафіксовано ціль передавання конфігурації \\
Сертифікат CA & Довірений сертифікат CA & Матеріал CA чинний і доступний & Зафіксовано ланцюжок довіри \\
Клієнтський сертифікат & CSR і випущений або імпортований сертифікат & Сертифікат узгоджений із CSR, пристроєм і профілем & Зафіксовано матеріал ідентифікації пристрою \\
Зв'язки & Користувач, пристрій, сертифікат, профіль і сервіс & Граф зв'язків узгоджений або виправлений & Зафіксовано цілісний доменний граф \\
Готовність матеріалів & Усі попередні матеріали та зв'язки & Перевірки довіри, авторизації й повноти успішні & Зафіксовано готовність до передавання \\
Підготовка та передавання конфігурації & Канонічна команда, ревізія і контрольна сума & Транзакцію створено і команда прийнята VPN або поставлена в чергу повтору & Зафіксовано стан «застосовано», «обмежено» або «помилка» \\
\bottomrule
\end{longtable}

Поточна реалізація повинна забезпечувати:
\begin{itemize}
  \item стійку чернетку із поточним кроком та вибраними доменними ідентифікаторами;
  \item умови переходів і заборону руху вперед без необхідних матеріалів;
  \item автоматичне застосування або ручне виправлення графа зв’язків;
  \item перевірку матеріалів сертифікатів, стан довіри та рішення про авторизацію;
  \item створення ідемпотентної транзакції підготовки VPN-конфігурації;
  \item доставку бажаного стану до VPN і відображення результату «застосовано», «обмежено» або «помилка»;
  \item повторні спроби, відновлення після перезапуску і звіряння стану після недоступності VPN.
\end{itemize}

\subsection{Імпорт OVPN}
Імпорт OVPN є допоміжним процесом первинного підключення. Система повинна аналізувати конфігурацію, визначати вбудовані блоки CA, сертифіката та ключа і параметри підключення, формувати попередній перегляд та створювати погоджені доменні об'єкти без збереження приватного ключа в IAS.

\subsection{CSR та сертифікати}
\begin{itemize}
  \item IAS повинен приймати зовнішній CSR, перевіряти формат, власника запиту та криптографічну узгодженість доступними засобами OTP.
  \item IAS повинен зв'язувати випущений сертифікат із CSR, пристроєм, користувачем і профілем безпеки.
  \item Перевірка строку чинності, видавця, відбитка, незмінності відкритого ключа та узгодженості зв’язків є обов'язковою.
  \item Перевірка CRL/OCSP у режимі реального часу належить до \textbf{цільового стану}, доки окрема інтеграція не реалізована і не протестована.
\end{itemize}

\subsection{Авторизація і довіра}
Поточна реалізація використовує профілі, політики, атрибути доменних об'єктів, граф зв’язків і довіру до сертифікатів для рішень у IAS/VPN сценарії. Це \textbf{часткова реалізація ABAC}: універсальна мова опису політик, зовнішні джерела атрибутів, геолокаційні правила та універсальний рушій ABAC належать до \textbf{цільового стану}.

\subsection{Передавання конфігурацій і звіряння стану}
\begin{itemize}
  \item Кожна керована зміна повинна мати канонічне представлення, ревізію і контрольну суму.
  \item Повторна доставка однакової команди повинна бути ідемпотентною.
  \item Застарілі ревізії не повинні перезаписувати новіший стан.
  \item Поточний і застосований стани повинні зберігатися після перезапуску.
  \item Осиротілий запис, невідповідність відбитка сертифіката, застаріла ревізія та неавторизований стан VPN повинні відображатися як інциденти звіряння стану.
  \item Дії з відновлення та виведення з експлуатації повинні бути аудованими та захищеними від повторного застосування старого бажаного стану.
\end{itemize}

\section{Нефункціональні вимоги}
\subsection{Архітектура і надійність}
\begin{itemize}
  \item Реалізація на Erlang/OTP із деревами нагляду і безпечними щодо перезапуску переходами станів.
  \item Авторитетне зберігання даних у KVS/Mnesia.
  \item Явні версії схем та логіка оновлення і міграції для стійких записів.
  \item Автоматизовані EUnit і Common Test сценарії для доменних сховищ, підготовки й передавання конфігурацій, RPC, звіряння стану та відновлення.
  \item Кількісні цільові показники SLA, місткості та затримки повинні підтверджуватися окремими навантажувальними та приймальними випробуваннями; до цього вони мають статус \textbf{Цільовий стан}.
\end{itemize}

\subsection{Безпека}
\begin{itemize}
  \item Приватні ключі користувачів і пристроїв не зберігаються в IAS.
  \item Секрети, ключовий матеріал і повні значення облікових даних не записуються в журнали аудиту та системні журнали.
  \item Захист транспортного рівня, автентифікація вузлів і розподіл секретів повинні бути визначені для кожного профілю розгортання.
  \item КСЗІ Г3, CAdES-X Long, ДСТУ 4145-2002, ДСТУ ГОСТ 28147:2009 та шифрування даних у сховищі належать до \textbf{нормативних вимог або цільового стану}, доки немає окремих реалізаційних і атестаційних доказів.
\end{itemize}

\subsection{Логування і аудит}
\begin{itemize}
  \item IAS повинен аудувати переходи життєвого циклу, випуск та імпорт, рішення про авторизацію, доставку конфігурацій, інциденти звіряння стану та дії відновлення.
  \item Записи аудиту повинні мати ідентифікатор об'єкта, ревізію, результат, причину та часову мітку.
  \item Сховище WORM та передавання до ELK/Loki і довгострокове нормативне зберігання є окремими вимогами до розгортання та нормативної відповідності, а не автоматичною властивістю поточної реалізації.
\end{itemize}

\subsection{Документація}
\begin{itemize}
  \item Технічні вимоги, техноробочий проєкт і політики повинні використовувати однакові статуси «Реалізовано», «Частково реалізовано», «Цільовий стан» і «Нормативна вимога».
  \item API/OpenAPI, настанови оператора та профілі безпеки створюються відповідно до фактично наявних інтерфейсів.
  \item BPMN/FSM-діаграми повинні спочатку відображати поточний життєвий цикл; виконання через BPE додається після реалізації відповідної інтеграції із середовищем виконання.
\end{itemize}

\section{Простежуваність вимог}
\begin{longtable}{p{2.3cm}p{5.4cm}p{2.2cm}p{4.1cm}}
\toprule
\textbf{ID} & \textbf{Вимога} & \textbf{Статус} & \textbf{Доказ / зона реалізації} \\
\midrule
IAS-FR-001 & Стійкий реєстр користувачів, пристроїв, сертифікатів, профілів і VPN-сервісів & Реалізовано & доменні модулі сховища, тести KVS/Mnesia \\
IAS-FR-002 & Перевірка CSR і X.509 & Реалізовано & модулі сертифікатів і налаштування, тести \\
IAS-FR-003 & Випуск сертифікатів через CA без зовнішнього командного інтерфейсу OpenSSL у середовищі виконання IAS & Цільовий стан & керований клієнт або адаптер CA \\
IAS-FR-004 & Майстер прив’язаного до пристрою налаштування VPN & Реалізовано & майстер, готовність матеріалів, зв’язки, авторизація \\
IAS-FR-005 & Доставка IAS $\rightarrow$ VPN через RPC & Реалізовано & модулі клієнта VPN і доставки, набори CT \\
IAS-FR-006 & Ревізія, контрольна сума, ідемпотентність і захист від застарілих даних & Реалізовано & стан доставки та підготовки конфігурацій, тести \\
IAS-FR-007 & Звіряння стану, відновлення осиротілих записів і виведення з експлуатації & Реалізовано & модулі звіряння й відновлення, набори CT \\
IAS-FR-008 & Керування процесами через BPE/BPMN & Цільовий стан & заплановане перенесення керування майстром \\
IAS-FR-009 & Адаптери CRL/OCSP, EST і CMC & Цільовий стан & відсутні у поточному середовищі виконання \\
IAS-NFR-001 & КСЗІ Г3, КЕП, криптографія за ДСТУ, шифрування даних у сховищі & Нормативна вимога & потребує окремої реалізації та підтвердних матеріалів \\
\bottomrule
\end{longtable}

\section{Додаток А. Межі поточної версії}
Поточна версія не повинна описуватися як така, що вже забезпечує повну універсальну ABAC-платформу, EST/CMC, перевірку CRL/OCSP у режимі реального часу, загальний набір інтерфейсів REST/OpenAPI, TLS 1.3 для всіх каналів, зберігання даних у RocksDB/NVMe, аудит у сховищі WORM, гарантовані 100000 пристроїв, 50 мс затримку, КСЗІ Г3, CAdES-X Long або реалізацію національної криптографії. Такі твердження можуть бути включені лише як «Цільовий стан» або «Нормативна вимога» з окремими критеріями приймання.

\section{Додаток Б. Цільова послідовність розвитку}
Послідовність робіт задається пріоритетами та залежностями, наведеними нижче. Наступний етап не повинен вважатися завершеним, доки не виконано його критерій результату.

\begin{longtable}{p{1.3cm}p{4.4cm}p{3.2cm}p{5.2cm}}
\toprule
\textbf{Пріоритет} & \textbf{Робота} & \textbf{Залежить від} & \textbf{Очікуваний результат} \\
\midrule
P0 & Синхронізувати \texttt{ias.tex} і \texttt{ias\_pro.tex} з поточним кодом & Немає & Документи коректно розмежовують реалізовані, частково реалізовані, цільові та нормативні вимоги \\
P1 & Повністю усунути запуск зовнішнього командного інтерфейсу OpenSSL із середовища виконання IAS, зберігши CMP або інший погоджений контракт через керований клієнт або адаптер CA & P0 & Випуск сертифікатів не залежить від локального виконуваного файла OpenSSL у IAS \\
P2 & Прогнати і стабілізувати повний життєвий цикл майстра налаштування VPN & P1 & Наскрізний сценарій від користувача і пристрою до передавання конфігурації відтворюється, перевіряється тестами та відновлюється після збоїв \\
P3 & Зафіксувати поточну FSM/BPMN-модель реального життєвого циклу & P2 & Стани, переходи, умови, помилки, повторні спроби та завершення описані без розбіжностей із кодом \\
P4 & Поетапно перенести керування процесом на \texttt{synrc/bpmn}/BPE, зберігаючи доменні сервісні модулі, перевірену поведінку та майстер як користувацький інтерфейс процесу & P3 & BPE є джерелом стану й керує переходами, а майстер залишається інтерфейсом над процесом: відображає його стан і надсилає команди \\
P5 & Повторно синхронізувати технічні вимоги, техноробочий проєкт і профілі безпеки & P4 & Документація, реалізація, випробування і профілі безпеки описують один погоджений стан системи \\
\bottomrule
\end{longtable}

Інвентаризація мережевих портів та сервісів у профілях безпеки виконується паралельно як окремий неблокувальний напрям, але її результати мають бути враховані не пізніше етапу P5.

\end{document}
